% SOURCE_STATEMENT: "1. It is known that there exists a good compactification of ALF spaces for L2-cohomology. What about its twistor space Q?(a) Does it compactify? (b) Does the complex structure extend over the compactification? 2. Study the Yang-Mills functional on open manifolds. For example, for manifolds which are conformally equivalent to a manifold with one cylindrical end, the energy of Yang-Mills instantons (if finite) is con-gruent mod Z to one of the Chern-Simons invariants of the boundary. What about different asymptotic geometries (ALF spaces, etc)?"
% Preparation and final self-audit: CONTRIBUTING.md prompts C and D; source/audit artifacts retained in the local physics_statements_2026_10_11 batch.
% !TeX program = lualatex
\documentclass[11pt]{article}
\usepackage[margin=25mm]{geometry}
\usepackage{fontspec}
\setmainfont{Latin Modern Roman}
\usepackage{amsmath,amssymb,mathtools}
\usepackage{unicode-math}
\setmathfont{Latin Modern Math}
\usepackage{xurl,hyperref}
\hypersetup{colorlinks=true,urlcolor=blue}
\setcounter{secnumdepth}{0}
\setlength{\parindent}{0pt}
\setlength{\parskip}{5pt}
\setlength{\emergencystretch}{3em}
\newcommand{\sref}[2]{\hyperlink{source:#1:#2}{[S#2]}}
\newcommand{\cataloguescope}[1]{\paragraph{Recorded target.}#1}
\newcommand{\R}{\mathbb R}
\newcommand{\C}{\mathbb C}
\newcommand{\N}{\mathbb N}
\newcommand{\Z}{\mathbb Z}
\newcommand{\Q}{\mathbb Q}
\newcommand{\F}{\mathbb F}
\newcommand{\eps}{\varepsilon}
\begin{document}
% UnsolvedMath ID 20002528; source catalogue code AIM-PHYSICS-0036
\setcounter{section}{20002527}
\section{A smoothness criterion for projective conic compactifications of ALF twistor models}\label{20002528}
{\small UnsolvedMath ID 20002528 · Mathematical Physics}
\subsection{Definitions and mathematical statement}

\paragraph{Shared notation.}$\mathbb N=\{1,2,\ldots\}$; the standard number fields have their usual meanings. All additional conventions are stated locally.


\paragraph{Definitions and notation.}A hyper-Kähler metric has three parallel complex structures satisfying the quaternionic relations. In four dimensions a gravitational instanton is a complete noncompact Ricci-flat geometry in the source’s self-dual setting. ALE and ALF describe different asymptotic geometries; ALF has an asymptotic circle direction over a three-dimensional base. The twistor space $Q$ of a self-dual or hyper-Kähler four-manifold is a complex threefold encoding its compatible complex structures. Compactification can mean smooth, complex-analytic or algebraic completion, with different requirements. The Yang--Mills functional is proportional to $\int|F_A|^2$, with normalization relevant modulo integers. Chern--Simons invariants depend on a connection on a boundary three-manifold and a fixed normalization.

\paragraph{Catalogue wording retained for traceability.}
For ALF gravitational instantons, determine whether the twistor space $Q$ admits a suitable compactification and whether its complex structure extends over that compactification. Study finite-energy Yang--Mills instantons on open manifolds with asymptotics such as ALF ends. Generalize the source’s cylindrical-end comparison between normalized instanton energy modulo $\mathbb Z$ and boundary Chern--Simons invariants, and decide whether the action remains quantized for other asymptotic geometries.
\paragraph{Scope and source qualification.}Both twistor and Yang–Mills questions are retained. The source warns that smoothness and irreducibility hypotheses matter for its Chern–Simons quantization statement; they must be checked, not dropped. A projective conic example does not settle arbitrary ALF twistor compactification. The catalogue title/status was paired with a machine-generated research attempt; that attempt is not peer-reviewed evidence of settlement. The statement covers the original source questions rather than only the title’s subcase.
\subsection{Short English statement}
Compactify ALF twistor spaces and determine Yang–Mills energy quantization for different noncompact asymptotics.
\subsection{Sources}
\begin{itemize}
\item[\hypertarget{source:20002528:1}{S1}] ULAM.AI and the UnsolvedMath Contributors, UnsolvedMath, record 20002528 (AIM-PHYSICS-0036). Catalogue curation under CC BY 4.0; definitions and English exposition added locally.
\url{https://huggingface.co/datasets/ulamai/UnsolvedMath}.
\item[\hypertarget{source:20002528:2}{S2}] American Institute of Mathematics, $L^2$ harmonic forms in geometry and string theory; notes by Anda Degeratu and Mark Haskins, March 2004, 21 March 2004, Etesi questions.
\url{https://aimath.org/WWN/l2harmonic/l2harmonic.pdf}.
\end{itemize}
\label{end:20002528}
\end{document}
